\section{Introduction}
\label{sec:intro}

Consider two tensor--train rank allocations with the same parameter count. The first has the smaller Frobenius reconstruction error for every compressed weight matrix, but removes a direction that is rarely activated and decisive on a safety capability. The second reconstructs the weights less accurately but preserves the teacher's conditional law on deployment prompts. A weight-space criterion necessarily selects the first allocation; a population functional criterion selects the second. This elementary possibility is consequential for large conditional generators because compression is ultimately judged through generated behavior, whereas tensor decompositions are usually fitted and compared through weight or activation reconstruction. Recent empirical work has documented this disconnect for post-training tensor compression~\citep{zagitov2026rethinking,swiftsvd2026}. A remaining statistical question is how to define the population target and intrinsic estimation price when one TT-compressed student must be learned from finitely many teacher queries.

Existing guarantees sharply constrain this question. Classical TT theory characterizes representation, approximation, and smooth fixed-rank geometry~\citep{oseledets2011tt,holtz2012manifolds}. Michel and Nouy's tree-tensor theory gives metric entropy and penalized empirical-risk model selection for general contrasts, with fast rates and near-minimax adaptation developed in bounded least squares~\citep{michel2022tree}. TT completion, recovery, regression, and quantile-regression theory already provides nonasymptotic rates, minimax lower bounds, or asymptotic inference in explicit observation models~\citep{imaizumi2017ttrank,zhou2022ttoi,qinwakin2024ttrecovery,qin2024ttregression,liu2024ttquantile,peng2025ttregression}. Matching low-rank adaptation bounds also extend to frozen nonlinear maps and locally quadratic cross-entropy settings~\citep{arunan2026lora}.

A second line of work concerns the observation channel and the criterion used for compression. Statistical analyses of knowledge distillation explain soft-label variance reduction and give generalization or optimization results in structured student models~\citep{phuong2019distillation,menon2021distillation,xie2026distillation}. Compression-based bounds control a network through its compressed description but do not analyze the present teacher-query experiment~\citep{arora2018compression,lotfi2024compression}. Fisher-, activation-, and KL-weighted matrix methods provide local task-sensitive objectives and calibration-aware rank allocation~\citep{fisherweighted2022,gfwsvd2025,iosvd2026,laser2026,rafft2026,statlora2026}; tensor-regression optimism likewise gives prediction-oriented rank selection for neural-network compression~\citep{shi2026optimism}. Among the primary sources located through 7 September 2026, we did not identify the particular oracle/lower-bound pair for the finite-query, teacher-relative conditional-density experiment on a TT weight quotient studied here. Moreover, a dimension-over-sample-size lower bound cannot be asserted without specifying what is observed from the teacher.

We study a finite-query distillation experiment. An observation consists of a deployment environment, a prompt, and one stochastic response from a fixed teacher conditional law. A candidate \(\lambda\) specifies tensorizations and exact, rank-separated TT strata for the compressed layers. The target is the weighted population conditional Kullback--Leibler divergence \(\cD_w(W,W^0)\), and the approximation error is the best joint risk of one \emph{common} deployable student. Our primary question is whether model selection can attain
\[
 \inf_{\lambda\in\Lambda}
 \left\{A_w^2(\lambda)+
 \frac{d_\lambda+\log H_\lambda+\log(1/\pi_\lambda)}{n}\right\},
 \qquad
 d_\lambda=\text{core parameters}-\text{TT gauge dimensions},
\]
and whether both terms are information-theoretically necessary. This experiment describes black-box or API-limited compression and finite-query distillation. We separately give soft-label and white-box counterexamples showing that the black-box lower bound does not transfer without additional assumptions.

Five issues make the question nonstandard. First, raw TT cores are nonidentifiable: adjacent invertible changes of basis leave the contracted tensor unchanged, so parameter counting overstates dimension. Second, quotient identifiability is not functional identifiability. Frozen network symmetries, dead activations, or unsupported prompts can make the conditional-law map singular even after the TT gauge has been removed. Third, approximation, model selection, and rank recovery interact badly near lower-rank strata, where a fixed-rank manifold loses reach. Fourth, capability certification reuses the prompts that selected ranks and fitted cores; ordinary holdout calibration does not justify this reuse. Finally, Fisher information supplies the local second-order geometry of correctly specified conditional KL\@. It is neither the global estimand nor the generic Hessian of an arbitrary downstream loss.

The analysis proceeds from the oracle theorem to its boundaries and extensions. On compact, separated TT quotient strata equipped with a finite pivot-gauge atlas, localized likelihood entropy is proportional to the gauge-adjusted dimension, and penalized likelihood gives a nonasymptotic oracle inequality. An Assouad hypercube in horizontal coordinates matches its estimation term; under local information-normal and off-chart separation conditions, a dense-teacher family also makes approximation and estimation jointly unavoidable for a fixed resource class. Classical local asymptotic minimax theory, instantiated on the quotient experiment, yields the constant \(\tfrac12\tr(G_w I_\pi^{-1})\). Separate access-model counterexamples show exact recovery from an identifiable logistic soft-label design and zero risk when a feasible teacher parameter is provided directly. Finally, we give a scoped functional-versus-Frobenius misordering result and apply algorithmic stability to simultaneous capabilities evaluated on the same independent prompt clusters used in fitting.

The paper makes five contributions.
\begin{enumerate}[leftmargin=*,label=(C\arabic*)]
    \item \textbf{Population TT oracle inequality.} Under a quantitative rank-separated quotient atlas and likelihood regularity conditions, we prove a finite-sample bound for one selected nonlinear TT student with joint population approximation error and complexity \(d_\lambda+\log H_\lambda+\log(1/\pi_\lambda)\), where \(H_\lambda\) is the atlas size; the TT dimension is calculated on the quotient, regularity factors are displayed, and optimization error is kept explicit.
    \item \textbf{Modelwise lower bounds and local constant.} We prove a \(d_\lambda/n\) lower bound on a regular TT stratum and, under explicit local information-normal and off-chart separation conditions, an \(A_w^2(\lambda)+d_\lambda/n\) lower bound for a fixed-budget teacher family. The quotient instantiation of the classical local asymptotic minimax constant becomes \(d_\lambda/2\) only under matched conditional-KL sampling.
    \item \textbf{Access and adaptation boundaries.} We give soft-label and white-box counterexamples to a universal \(d_\lambda/n\) lower bound, and show that uniform rank recovery is impossible within an \(n^{-1/2}\) functional neighborhood of a lower-rank stratum. We state conditional underfit-gap and overfit-penalty conditions sufficient for consistent penalized rank selection.
    \item \textbf{Functional improvement with correct scope.} We quantify when an anisotropic output metric reverses the ordering induced by Frobenius error and show that functional validation selects the better candidate exponentially fast in its risk gap. This is a strict improvement on a constructed class, not a universal domination claim.
    \item \textbf{Application-level same-data certification.} Adapting classical uniform-stability concentration to a fixed coded family, we give simultaneous capability upper bands valid after fitting and selection, along with a least-cost feasibility guarantee. This secondary corollary does not replace protected-holdout testing when stability is unavailable.
\end{enumerate}

The results are stated for finite or length-truncated outputs and for a concrete nonlinear frozen-prefix/frozen-head block. Extending the constants to a complete transformer requires restricted-Jacobian conditioning, sequence-level tail control, multi-layer chart geometry, and a certified optimization guarantee. We expose these requirements rather than treating them as consequences of parameter count.
